% ---- pg-249 （书页 237） ----

\phantomsection\label{eq:6-49}
\begin{equation}
\mathbf{S}_p=\frac{1}{\sum\limits_{l}(n_l-1)}\left\{(n_1-1)\mathbf{S}_1+(n_2-1)\mathbf{S}_2+\cdots+(n_g-1)\mathbf{S}_g\right\} \tag{6-49}
\end{equation}

博克斯检验基于 $-2\ln\Lambda$ 的样本分布的 $\chi^2$ 近似（参见结论 5.2）．令 $-2\ln\Lambda=M$（博克斯 $M$ 统计量），得到
\phantomsection\label{eq:6-50}
\begin{equation}
M=\left[\sum\limits_{l}(n_l-1)\right]\ln|\mathbf{S}_p|-\sum\limits_{l}\left[(n_l-1)\ln|\mathbf{S}_l|\right] \tag{6-50}
\end{equation}

如果零假设是真的，各个体的样本协方差矩阵不会差太多，因此，与融合协方差矩阵也不会差太多．在这种情况下，\hyperref[eq:6-48]{式(6-48)}中的行列式之比都接近 1，$\Lambda$ 也将接近 1，并且博克斯 $M$ 统计量会比较小．如果零假设是假的，样本协方差矩阵可能差很多，而且它们的行列式的差异会更加明显．在这种情况下，$\Lambda$ 比较小，$M$ 相对比较大．为了说得更清楚，注意到融合协方差矩阵的行列式 $|\mathbf{S}_p|$，会处在这单组协方差矩阵行列式 $|\mathbf{S}_l|$ 的“中间”的某个位置附近．随着后者的行列式的值越不同，\hyperref[eq:6-48]{式(6-48)}中的比率的积将越接近 0．事实上，随着 $|\mathbf{S}_l|$ 的增加，$|\mathbf{S}_{(1)}|/|\mathbf{S}_p|$ 降低乘积的效果比 $|\mathbf{S}_{(g)}|/|\mathbf{S}_p|$ 增加乘积的效果更大，这里 $|\mathbf{S}_{(1)}|$ 和 $|\mathbf{S}_{(g)}|$ 分别为这些行列式中的最小值和最大值．

\begin{mvframe}[协方差矩阵相等性的博克斯检验]
令
\phantomsection\label{eq:6-51}
\begin{equation}
u=\left[\sum\limits_{l}\frac{1}{(n_l-1)}-\frac{1}{\sum\limits_{l}(n_l-1)}\right]\left[\frac{2p^2+3p-1}{6(p+1)(g-1)}\right] \tag{6-51}
\end{equation}
这里，$p$ 是变量的数目，$g$ 是群组的数目．则
\phantomsection\label{eq:6-52}
\begin{equation}
C=(1-u)M=(1-u)\left\{\left[\sum\limits_{l}(n_l-1)\right]\ln|\mathbf{S}_p|-\sum\limits_{l}\left[(n_l-1)\ln|\mathbf{S}_l|\right]\right\} \tag{6-52}
\end{equation}
近似服从 $\chi^2$ 分布，其自由度为
\phantomsection\label{eq:6-53}
\begin{equation}
\nu=g\frac{1}{2}p(p+1)-\frac{1}{2}p(p+1)=\frac{1}{2}p(p+1)(g-1) \tag{6-53}
\end{equation}
在置信水平 $\alpha$，如果 $C>\chi_{p(p+1)(g-1)/2}^{2}(\alpha)$，则拒绝零假设 $H_0$．
\end{mvframe}

当 $n_l$ 超过 20 并且 $p$ 和 $g$ 都不超过 5 时，博克斯的 $\chi^2$ 近似效果不错．在上述条件不成立的情况下，对于 $M$ 的样本分布，博克斯（参见文献[7]，[8]）提供了一个更加精确的 $F$ 分布近似．

\noindent\textbf{例 6.12}\quad （协方差矩阵相等性检验——小型疗养所）

我们在例 6.10 中介绍了威斯康星州小型疗养所的数据．我们列出了 $p=4$ 和 $g=3$ 的样本协方差矩阵．假设数据是多元正态的，我们检验零假设
\[
H_0:\boldsymbol{\Sigma}_1=\boldsymbol{\Sigma}_2=\boldsymbol{\Sigma}_3=\boldsymbol{\Sigma}
\]
利用例 6.10 中的信息，我们有 $n_1=271$，$n_2=138$，$n_3=107$ 和 $|\mathbf{S}_1|=2.783\times10^{-8}$，$|\mathbf{S}_2|=89.539\times10^{-8}$，$|\mathbf{S}_3|=14.579\times10^{-8}$ 以及 $|\mathbf{S}_p|=17.398\times10^{-8}$．对行列式的值取自然对数，有
\[
\ln|\mathbf{S}_1|=-17.397,\quad \ln|\mathbf{S}_2|=-13.926,\quad \ln|\mathbf{S}_3|=-15.741,\quad \ln|\mathbf{S}_p|=-15.564
\]
我们可以计算得
\[
u=\left(\frac{1}{270}+\frac{1}{137}+\frac{1}{106}-\frac{1}{270+137+106}\right)\left[\frac{2\times4^2+3\times4-1}{6\times(4+1)(3-1)}\right]=0.013\,3
\]
